\section{Conclusion}

We formulated context selection as a risk-controlled decision problem in which
the router chooses not only which auxiliary context to use but whether to use
any. A harmful-selection budget turns the reliability of a utility model into
a coverage number: conformal risk control calibrates a single threshold from
deployment states, the expected number of harmful selections per decision
stays within the budget, and every decision that cannot be certified becomes
an abstention rather than a risky selection. The same score supports a second
objective, a bound on the harmful share of accepted decisions, and the two
operating points bracket the trade-off an operator faces. Interval
certification, the natural alternative, retains four to nine times fewer
decisions at the same nominal level because its radius is tied to the
estimation error rather than to the accepted set. Thresholds calibrated on
some benchmarks transfer to another with a shift margin smaller than the
budget at strict levels, and the margin grows as the budget loosens. Reliable
knowledge-augmented prediction therefore requires an explicit decision about
when to abstain, and that decision can be calibrated, reported, and audited.
